\documentclass[10pt]{article}
\usepackage[margin=0.78in]{geometry}
\usepackage{booktabs}
\usepackage{graphicx}
\usepackage{microtype}
\usepackage{caption}
\title{Frozen-Encoder Spectral Transfer: SAC versus Actor+Critic JEPA}
\author{}
\date{September 23, 2026 --- RTX 3090}
\begin{document}
\maketitle

\section*{Question}
Does the JEPA-pretrained representation improve adaptation when its visual
encoder is held fixed after a terrain spectrum shift?  This focused comparison
keeps only the two relevant strategies: vanilla SAC and SAC with independent
actor and critic JEPA objectives.  The CDD/ConvNeXt experiment is intentionally
excluded: its preprocessing/runtime characteristics are a separate question.

\section*{Protocol}
Both source agents were trained for 50k transitions at spectral exponent
$\alpha=3$.  Each was initialized from its own endpoint and adapted for a
fresh 20k transitions at $\alpha=2$ and $\alpha=4$.  During each transfer, the
actor CNN/encoder and both independent critic CNN/encoders were frozen.  The
policy head, Q heads, entropy temperature, and (for JEPA) predictor heads
remained trainable.  Optimizer, replay buffer, and environments were fresh.
Numbers are deterministic held-out mean speed $\pm$ standard deviation across
four held-out terrain seeds; this is one training seed per condition.

\begin{figure}[h]
\centering
\includegraphics[width=.96\linewidth]{output/pdf/frozen_encoder_adaptation_curve.png}
\caption{The matched frozen-encoder adaptation curves.  JEPA is clearly ahead
at 10k on both shifts; the 20k endpoints converge.}
\end{figure}

\begin{center}\small
\begin{tabular}{llrr}
\toprule
Target $\alpha$ & Method & 10k transfer & 20k transfer \\
\midrule
2 & Vanilla SAC & $0.315\pm0.125$ & \textbf{$2.917\pm0.022$} \\
2 & Actor+critic JEPA & \textbf{$1.097\pm0.065$} & $2.842\pm0.039$ \\
\addlinespace
4 & Vanilla SAC & $2.309\pm0.097$ & $3.050\pm0.075$ \\
4 & Actor+critic JEPA & \textbf{$2.510\pm0.028$} & \textbf{$3.090\pm0.052$} \\
\bottomrule
\end{tabular}
\end{center}

\section*{Result}
JEPA improves early adaptation under fixed encoders: at 10k it is higher by
$0.782$ speed at $\alpha=2$ and $0.201$ at $\alpha=4$.  By 20k, the advantage
is not consistent: vanilla SAC is higher by $0.075$ at $\alpha=2$, while
actor+critic JEPA is higher by $0.040$ at $\alpha=4$.  Thus the defensible
claim is an early adaptation benefit, not a proven final-performance gain.

\section*{Verification artifacts}
Each finalized transfer directory contains a full analysis manifest, learning
curve, held-out trajectory/UMAP, episode diagnostics, pairwise UMAP, and CKA.
The four run directories are listed below:
\begin{center}\small\ttfamily
dumps/sac\_frozen\_alpha2\_transfer/\\
dumps/sac\_frozen\_alpha4\_transfer/\\
dumps/sac\_jepa\_actorcritic\_frozen\_alpha2\_transfer/\\
dumps/sac\_jepa\_actorcritic\_frozen\_alpha4\_transfer/
\end{center}
The CKA matrices are identically 1.00 across the two transfer checkpoints, as
expected when the encoders are frozen.

\end{document}
